\chapter{Grundlagen}
\label{chap:theory}

\offen{Die fachlichen Grundlagen mit zugänglichen Originalquellen belegen.}
Als reines Satzbeispiel wird in \eqnref{eq:ohm} ein ohmsches Modell bei
konstanter Temperatur verwendet. Ob dieses Modell für einen konkreten
Versuch geeignet ist, wird anhand des Bauteils und der Versuchsbedingungen
geprüft.
\begin{equation}
  R = \frac{U}{I}
  \label{eq:ohm}
\end{equation}
Dabei bezeichnet $R$ den elektrischen Widerstand in \unit{\ohm}, $U$ die
Spannung in \unit{\volt} und $I$ die Stromstärke in \unit{\ampere}.
Die Formelzeichen werden zusätzlich in der Übersicht am Ende aufgeführt.

\section{Quellen und Zitate}
Hier wird nur die Position eines Quellenverweises demonstriert; der
folgende Beleg ist ausdrücklich ein zu ersetzender Platzhalter.\mcicite{quelle-ersetzen}
Bei einer echten fachlichen Übernahme wird der Beleg direkt an die
übernommene Aussage gesetzt. Ein kurzes wörtliches Zitat wird beispielsweise
als \mciquote{OFFEN: exakten Originalwortlaut einsetzen} gesetzt; die
Originalquelle und die genaue Fundstelle werden ergänzt.
